% =========================================================
% Revised section opening (replaces the old 030C introduction)
% =========================================================
\section{Illustrative Stakeholder-Specific Narratives}
\label{SEC:stakeholder_narrative_examples}

To demonstrate the effect of stakeholder-specific evidence selection, six
narratives were constructed around the same analytical instance: Kingston upon
Hull 029E (\texttt{E01033104}), the city-centre and Old Town LSOA, with the
GeoSensitivity map view active. The planner narrative adds a contrasting
low-sensitivity case, Kingston upon Hull 003C (\texttt{E01012877}) in Orchard
Park. The underlying structured evidence is therefore held constant while the
selection, prioritisation, and communication of that evidence changes
according to stakeholder needs.

All values come from the final Hull design: one XGBoost model explained by
the Feature branch (30 non-spatial predictors) and by the area-data Geo branch,
in which masking an LSOA replaces all of its inputs with training medians
($N=4{,}000$ perturbations, seed 42). The examples were written manually
according to the dashboard prompt and are therefore \emph{illustrative
prompt-conforming narratives rather than API outputs}. Place names and
streets come from the \texttt{place\_context} supplied to the prompt, built
from OpenStreetMap features inside the official 2021 LSOA boundaries.

The examples are intended to demonstrate differences in evidence selection,
rather than merely differences in vocabulary. Each subsection therefore
presents the same analytical instance through a different stakeholder lens,
with the dashboard evidence on the left and the narrative on the right.

% ---------------------------------------------------------
% [Resident / Public and City Planner subsections: unchanged]
% ---------------------------------------------------------


% =========================================================
\subsection{Police / Community Safety}
% =========================================================

\begin{tcolorbox}[
    enhanced,
    colback=white,
    colframe=policecol,
    boxrule=0.6pt,
    arc=2mm,
    borderline west={3pt}{0pt}{policecol},
    left=2.5mm,
    right=2.5mm,
    top=2.5mm,
    bottom=2.5mm,
    before skip=4pt,
    after skip=4pt
]

{\small
\textbf{Stakeholder question.}
\emph{Which local signals dominate the explanation here, is the pattern
persistent or recent, and does the area warrant contextual review?}
}

\vspace{2mm}

\noindent
\begin{minipage}[t]{0.41\linewidth}
\vspace{0pt}

{\scriptsize\bfseries DASHBOARD EVIDENCE}

\vspace{0.8mm}

\centering
\includegraphics[
    width=\linewidth,
    height=0.26\textheight,
    keepaspectratio
]{Chapter6/figs/police_E01033104_rm6_reliance_map.png}

\vspace{0.3mm}

\raggedright
{\scriptsize\color{black!60}
\textbf{Local feature view.}
Reliance on the \textbf{six-month crime average} is highest in 029E
(\textbf{48.5; rank 1 of 168}) and is spatially clustered
(Moran's $I=\textbf{0.45}$, $p=0.001$).
}

\vspace{1.5mm}

\centering
\includegraphics[
    width=\linewidth,
    height=0.17\textheight,
    keepaspectratio
]{Chapter6/figs/police_E01033104_crime_timeline.png}

\vspace{0.3mm}

\raggedright
{\scriptsize\color{black!60}
\textbf{Persistence view.}
Recorded crime has stayed between \textbf{200 and 226} a month in 2026,
against a Hull median of about \textbf{14}.
}

\end{minipage}\hfill%
\begin{minipage}[t]{0.55\linewidth}
\vspace{0pt}

{\scriptsize\bfseries STAKEHOLDER-SPECIFIC NARRATIVE}

\vspace{0.8mm}

\footnotesize
\RaggedRight

Across Hull, the model leans on recent crime history. The six-month crime
average is the strongest local signal in \textbf{144 of 168 LSOAs}, and areas
that rely on it heavily tend to sit next to each other.

\vspace{1.2mm}

Kingston upon Hull 029E (\texttt{E01033104}) covers the city centre and Old
Town around Ferensway, Paragon Square, Carr Lane, Market Place and Castle
Street. It relies on that signal more than any other area
(\textbf{rank 1 of 168}). Its top three local signals are the six-month
average (\textbf{48.5}), the three-month average (\textbf{37.8}) and the count
from two months earlier (\textbf{16.4}); crime history makes up
\textbf{83\%} of its local explanation, against a Hull median of
\textbf{53\%}.

\vspace{1.2mm}

The pattern is \textbf{persistent rather than a recent spike}. The area
recorded \textbf{217.5} crimes a month on average in January--April 2026,
\textbf{7.1\%} of all recorded crime in Hull, with every month between
\textbf{200 and 226} and a six-month average of \textbf{214.5} in April.

\vspace{1.2mm}

The model is also more sensitive to this area's data than to any other
(\textbf{rank 1 of 168}). 029E sits in a group of \textbf{nine} areas with
unusually high sensitivity, running through the Old Town, west along Hessle
Road and Anlaby Road into Dairycoates around Coltman Street, and north-west
along Spring Bank into Stepney around Argyle Street.

\vspace{1.2mm}

The area is \textbf{worth closer contextual review}; this evidence describes
how the model uses the data and is \textbf{not a basis for deployment}.

\end{minipage}

\vspace{2mm}
\hrule
\vspace{1mm}

{\scriptsize
\textbf{Stakeholder emphasis.}
The police narrative prioritises the \textbf{dominant local crime-history
signals} and adds \textbf{persistence evidence} that the other narratives
omit. It keeps the spatial context brief and closes by preventing the
explanation from being read as a deployment recommendation.
}

\end{tcolorbox}


% =========================================================
\subsection{Manager / Strategic Decision Maker}
% =========================================================

\begin{tcolorbox}[
    enhanced,
    colback=white,
    colframe=managercol,
    boxrule=0.6pt,
    arc=2mm,
    borderline west={3pt}{0pt}{managercol},
    left=2.5mm,
    right=2.5mm,
    top=2.5mm,
    bottom=2.5mm,
    before skip=4pt,
    after skip=4pt
]

{\small
\textbf{Stakeholder question.}
\emph{What are the main city-wide findings, what is exceptional, can the
results be trusted, and where should analytical review be concentrated?}
}

\vspace{2mm}

\noindent
\begin{minipage}[t]{0.41\linewidth}
\vspace{0pt}

{\scriptsize\bfseries DASHBOARD EVIDENCE}

\vspace{0.8mm}

\centering
\includegraphics[
    width=\linewidth,
    height=0.17\textheight,
    keepaspectratio
]{Chapter6/figs/manager_geo_ranked.png}

\vspace{0.3mm}

\raggedright
{\scriptsize\color{black!60}
\textbf{Concentration view.}
One area stands out (\textbf{22.33}, about three times the next); most
areas sit just below zero.
}

\vspace{1.5mm}

\centering
\includegraphics[
    width=\linewidth,
    height=0.25\textheight,
    keepaspectratio
]{Chapter6/figs/manager_geo_top20.png}

\vspace{0.3mm}

\raggedright
{\scriptsize\color{black!60}
\textbf{Expert agreement.}
\textbf{18 of the 20} most sensitive areas lie in the regions the expert
highlighted.
}

\end{minipage}\hfill%
\begin{minipage}[t]{0.55\linewidth}
\vspace{0pt}

{\scriptsize\bfseries STAKEHOLDER-SPECIFIC NARRATIVE}

\vspace{0.8mm}

\footnotesize
\RaggedRight

\textbf{Main finding.} The model relies overwhelmingly on recent crime
history: the six-month crime average is the leading local factor in
\textbf{144 of 168 LSOAs}. Social and built-environment factors play a
secondary role.

\vspace{1.2mm}

\textbf{What is exceptional.} Geographic sensitivity is concentrated. Only
\textbf{48 of 168} areas are above the Hull average, and the 20 most
sensitive account for \textbf{46\%} of the total effect of replacing each
area's data. The city centre and Old Town, Kingston upon Hull 029E
(\texttt{E01033104}), is the single exceptional case at \textbf{22.33},
about \textbf{three times} the next area, 030B in Dairycoates
(\textbf{7.18}). It also accounts for \textbf{7.1\%} of Hull's recorded
crime.

\vspace{1.2mm}

\textbf{Can it be trusted?} The results line up with local expertise:
\textbf{18 of the 20} most sensitive areas fall in the three regions the
expert highlighted (\textbf{90\%}, against about \textbf{35\%} by chance).
The rankings are unchanged when neighbouring areas are added in any
direction, and replacing the data of the top 20 areas moves predictions far
more than replacing random areas. Across the ten candidate models, however,
more accurate models did not give closer agreement with the expert, so
model accuracy alone is not a guarantee of explanation quality.

\vspace{1.2mm}

\textbf{Next step.} Review the evidence for the city-centre cluster, and for
029E in particular, before relying on the geographic rankings elsewhere.
None of these findings is causal.

\end{minipage}

\vspace{2mm}
\hrule
\vspace{1mm}

{\scriptsize
\textbf{Stakeholder emphasis.}
The manager receives a \textbf{strategic synthesis} organised as finding,
exception, confidence and action. It reports whether the pattern is
\textbf{typical or exceptional} and how far the evidence can be trusted, and
it recommends an \textbf{evidence-review action} rather than an operational
intervention.
}

\end{tcolorbox}


% =========================================================
\subsection{GeoAI Researcher / Model Analyst}
% =========================================================

\begin{tcolorbox}[
    enhanced,
    colback=white,
    colframe=researchercol,
    boxrule=0.6pt,
    arc=2mm,
    borderline west={3pt}{0pt}{researchercol},
    left=2.5mm,
    right=2.5mm,
    top=2.5mm,
    bottom=2.5mm,
    before skip=4pt,
    after skip=4pt
]

{\small
\textbf{Stakeholder question.}
\emph{How is Geo-SMILE behaving technically for this instance, and how
faithful, spatially coherent, and reproducible is the explanation?}
}

\vspace{2mm}

\noindent
\begin{minipage}[t]{0.41\linewidth}
\vspace{0pt}

{\scriptsize\bfseries DASHBOARD EVIDENCE}

\vspace{0.8mm}

\centering
\includegraphics[
    width=\linewidth,
    height=0.2\textheight,
    keepaspectratio
]{Chapter6/figs/researcher_beta_vs_exact_shift.png}

\vspace{0.3mm}

\raggedright
{\scriptsize\color{black!60}
\textbf{Surrogate check.}
Signed $\beta$ tracks the exact single-LSOA $W_1$ shift
($\rho_s=\textbf{0.982}$); 029E is the extreme point.
}

\vspace{1.5mm}

\centering
\includegraphics[
    width=\linewidth,
    height=0.23\textheight,
    keepaspectratio
]{Chapter6/figs/researcher_lisa_map.png}

\vspace{0.3mm}

\raggedright
{\scriptsize\color{black!60}
\textbf{Spatial structure.}
029E is a High--High member (local $I=\textbf{12.08}$, $p=0.003$) of the
single nine-LSOA HH cluster.
}

\end{minipage}\hfill%
\begin{minipage}[t]{0.55\linewidth}
\vspace{0pt}

{\scriptsize\bfseries STAKEHOLDER-SPECIFIC NARRATIVE}

\vspace{0.8mm}

\footnotesize
\RaggedRight

Kingston upon Hull 029E (\texttt{E01033104}) has a signed area-data
Geo-SMILE coefficient of $\beta=\textbf{22.33}$ (intercept \textbf{9.79}),
\textbf{rank 1/168}. Its single-LSOA Wasserstein prediction shift is
\textbf{176.7} crimes per month, against a median of \textbf{6.4} and a 90th
percentile of \textbf{22.0}. The signed ranking tracks this exact shift
($\rho_s=\textbf{0.982}$); the $|\beta|$ ranking does not
($\rho_s=\textbf{$-0.29$}$), so the sign marks a relative contribution to the
coalition $W_1$, not a direction.

\vspace{1.2mm}

Global Moran's $I$ on $\beta$ is \textbf{0.22} ($p=\textbf{0.001}$, queen
contiguity, \textbf{999 permutations}); \textbf{142 of 168} LSOAs are
non-significant. 029E is HH (local $I=\textbf{12.08}$, unadjusted
$p=\textbf{0.003}$) in one contiguous HH cluster of \textbf{nine}; its Old
Town neighbours 029C, 029H, 029F and 029G rank \textbf{4, 12, 16, 17}.

\vspace{1.2mm}

Locally, \texttt{crime\_rolling\_mean\_6} (\textbf{$-48.52$}) and
\texttt{crime\_rolling\_mean\_3} (\textbf{$-37.80$}) lead, matching the top
two of the global Feature ranking ($|\beta|=\textbf{5.14}$, \textbf{1.25});
total local $|\beta|$ is \textbf{152.4} against a Hull median of
\textbf{9.8}. A separate joint location player (coordinates added to the
model) contributes only \textbf{$-0.13$} here and ranks \textbf{26/31}
globally, so the Geo signal does not reduce to position.

\vspace{1.2mm}

Held-out surrogate $R^2$ is \textbf{0.977} (Feature) and \textbf{0.948}
(Geo). Geo top-20 ATT accuracy is \textbf{18/20} (AUROC \textbf{0.701}),
top-20 minus random $\Delta W_1=\textbf{26.22}$ ($p=\textbf{0.001}$), and
cardinal-addition Jaccard is \textbf{1.0}. Across five seeds 029E has SD
\textbf{0.65} (\textbf{2.8\%} of its mean); Geo rank Spearman is
\textbf{0.963}.

\end{minipage}

\vspace{2mm}
\hrule
\vspace{1mm}

{\scriptsize
\textbf{Stakeholder emphasis.}
This version retains \textbf{exact values}, the sign convention of each
branch, spatial statistics, the location-player contrast and the
\textbf{explanation-quality diagnostics}. Its purpose is technical
interrogation of the framework rather than rapid decision support.
}

\end{tcolorbox}


% =========================================================
\subsection{Governance / Assurance}
% =========================================================

\begin{tcolorbox}[
    enhanced,
    colback=white,
    colframe=governancecol,
    boxrule=0.6pt,
    arc=2mm,
    borderline west={3pt}{0pt}{governancecol},
    left=2.5mm,
    right=2.5mm,
    top=2.5mm,
    bottom=2.5mm,
    before skip=4pt,
    after skip=4pt
]

{\small
\textbf{Stakeholder question.}
\emph{Is the explanation traceable and reproducible, and which limitations
require explicit review before the evidence is relied upon?}
}

\vspace{2mm}

\noindent
\begin{minipage}[t]{0.41\linewidth}
\vspace{0pt}

{\scriptsize\bfseries DASHBOARD EVIDENCE}

\vspace{0.8mm}

\centering
\includegraphics[
    width=\linewidth,
    height=0.18\textheight,
    keepaspectratio
]{Chapter6/figs/governance_seed_consistency.png}

\vspace{0.3mm}

\raggedright
{\scriptsize\color{black!60}
\textbf{Reproducibility.}
Five independently seeded runs: 029E varies by \textbf{2.8\%}; lower
ranks vary more.
}

\vspace{1.5mm}

\centering
\includegraphics[
    width=\linewidth,
    height=0.24\textheight,
    keepaspectratio
]{Chapter6/figs/governance_stability_four_directions.png}

\vspace{0.3mm}

\raggedright
{\scriptsize\color{black!60}
\textbf{Stability.}
The top 20 is unchanged under all four neighbouring-LSOA additions
(Jaccard \textbf{1.0}).
}

\end{minipage}\hfill%
\begin{minipage}[t]{0.55\linewidth}
\vspace{0pt}

{\scriptsize\bfseries STAKEHOLDER-SPECIFIC NARRATIVE}

\vspace{0.8mm}

\footnotesize
\RaggedRight

\textbf{Provenance.} Kingston upon Hull 029E (\texttt{E01033104}), the city
centre and Old Town, has a GeoSensitivity of \textbf{22.33}
(\textbf{rank 1 of 168}). It was produced by the area-data Geo branch: one
XGBoost model, \textbf{4,000} perturbations with seed \textbf{42}, each
replacing all inputs of 5--10 LSOAs with training medians. The expert
regions and the rule assigning LSOAs to them were fixed before the ranking
was produced.

\vspace{1.2mm}

\textbf{Reproducibility.} Across five seeds 029E's coefficient varies by
\textbf{2.8\%} (SD \textbf{0.65}), and the ranking is stable overall (rank
Spearman \textbf{0.963}). Membership of the top fifth is less stable for Geo
than for Feature (Jaccard \textbf{0.77} versus \textbf{1.00}), so ranks near
that boundary should not be over-read. Adding a neighbouring area in any
direction leaves the top 20 unchanged.

\vspace{1.2mm}

\textbf{Limitations requiring review.}
\begin{itemize}[leftmargin=*, itemsep=0pt, topsep=1pt]
\item The coefficient sign is \textbf{not a direction} of predicted crime.
\item The expert reference is a \textbf{single rater}; LISA $p$-values are
\textbf{not adjusted} for multiple comparisons.
\item Across ten models, \textbf{ATT Faithfulness was not demonstrated}
($r=\textbf{0.15}$ Geo, \textbf{$-0.12$} Feature).
\item 029E is about \textbf{three times} the next area, so it dominates the
geographic picture and warrants separate model review.
\end{itemize}

\vspace{0.8mm}

None of these results is causal.

\end{minipage}

\vspace{2mm}
\hrule
\vspace{1mm}

{\scriptsize
\textbf{Stakeholder emphasis.}
Although the governance narrative uses technical metrics, its purpose differs
from the researcher's. It foregrounds \textbf{provenance},
\textbf{reproducibility} and \textbf{limitations}, including the negative
ATT Faithfulness result, rather than a comprehensive technical account of
the selected case.
}

\end{tcolorbox}
